%% =============================================================================
%% APPENDICES - content for all appendix sections
%% =============================================================================
%% Note: The \appendix command and TOC are in discussion.tex.
%% This file contains the actual appendix sections.

\section{Best practices and coverage gaps} \label{app:best-practices}
\setcounter{table}{0}
\setcounter{figure}{0}

This appendix provides practical guidance for effective use of
\pkg{prompt2analytics} and documents method coverage gaps.

\subsection{Pre-analysis checklist}

Before running an analysis through \pkg{prompt2analytics}, users should verify:
\begin{enumerate}
  \item \emph{Method availability}: Confirm the required method is implemented.
    Table~\ref{tab:method-coverage} documents coverage gaps;
    Appendix~\ref{app:implementation} provides method-specific caveats.
  \item \emph{Data requirements}: Verify the dataset meets method assumptions.
    Panel methods require entity and time identifiers; instrumental
    variables (IV) require instruments; difference-in-differences (DiD)
    requires treatment and time indicators.
  \item \emph{Variable naming}: Use clear, unambiguous column names. The LLM
    extracts variable names from natural language; names like \code{x1} or
    \code{var\_1} are more error-prone than \code{wage} or \code{education}.
  \item \emph{Specification clarity}: Specify methods explicitly when
    precision matters. ``Estimate a within-transformation fixed effects model
    with clustered standard errors at the firm level'' yields more reliable
    results than ``run a regression with firm effects.''
\end{enumerate}

\subsection{Recommended workflows}

The system is most effective for:
\begin{itemize}
  \item \emph{Exploratory analysis}: Quickly testing specifications before
    committing to a formal analysis pipeline.
  \item \emph{Teaching and learning}: Students can explore econometric methods
    conversationally while seeing which tools and parameters are invoked.
  \item \emph{Rapid prototyping}: Researchers can iterate on model specifications
    in natural language before writing production code.
  \item \emph{Verification}: Cross-checking results from other software by
    running the same analysis through an independent implementation.
\end{itemize}

\noindent Users should prefer \proglang{R}, \proglang{Stata}, or \proglang{Python}
when methods are unavailable (Bayesian inference, Heckman selection), maximum
performance is critical (\pkg{fixest} outperforms for high-dimensional fixed effects (HDFE) at scale),
reproducibility requires code scripts, or publication-quality visualizations
are needed.

\subsection{Method coverage gaps}

While method coverage is broad (over 200 methods), several important categories
remain unimplemented. Table~\ref{tab:method-coverage} summarizes coverage
gaps and recommends alternative tools.

\begin{table}[htbp]
\centering
\begin{threeparttable}
\caption{Method coverage gaps}
\label{tab:method-coverage}
\small
\begin{tabular}{p{3.5cm}p{5cm}p{4cm}}
\toprule
Method Category & Status & Recommended Alternative \\
\midrule
Bayesian inference & Not implemented & \proglang{R}: \pkg{rstanarm}, \pkg{brms}; \proglang{Python}: \pkg{PyMC} \\
Heckman selection & Not implemented & \proglang{R}: \pkg{sampleSelection}; \proglang{Stata}: \code{heckman} \\
Conley spatial SEs & Not implemented & \proglang{R}: \pkg{fixest}, \pkg{lfe} \\
Full SARIMA (arbitrary seasonal) & Partial support & \proglang{R}: \pkg{forecast} \\
Quantile IV & Not implemented & \proglang{R}: \pkg{quantreg} \\
Structural equation modeling & Not implemented & \proglang{R}: \pkg{lavaan} \\
\addlinespace
Weak instrument inference & $F$-stat only; no Stock-Yogo CVs & \proglang{R}: \pkg{ivmodel}, \pkg{AER} \\
\bottomrule
\end{tabular}
\begin{tablenotes}[flushleft]
\item \emph{Note:} Lists unimplemented or partially implemented method
categories with recommended alternatives in \proglang{R}, \proglang{Stata},
or \proglang{Python}.
\end{tablenotes}
\end{threeparttable}
\end{table}

\noindent Users encountering these gaps should consult the listed
alternatives; future releases aim to address Bayesian inference and
Heckman selection as priorities.

\section{Deployment guide} \label{app:deployment}
\setcounter{table}{0}
\setcounter{figure}{0}

Full platform-specific installation instructions (Linux, macOS, Windows/WSL2,
Docker) are maintained in the repository's \code{README.md} at
\url{https://github.com/umatter/prompt2analytics}. Pre-compiled binaries
for all major platforms are available on the GitHub releases page.
This appendix documents only the MCP client configuration and hardware
requirements referenced from Section~\ref{sec:deployment}.

\subsection{MCP client configuration}

For Claude Desktop, add to \code{claude\_desktop\_config.json}:

\begin{verbatim}
{
  "mcpServers": {
    "prompt2analytics": {
      "command": "/path/to/p2a-mcp",
      "args": ["--transport", "stdio"]
    }
  }
}
\end{verbatim}

\noindent Other MCP-compatible clients (Cursor, Windsurf, custom agents)
use the same \code{--transport stdio} interface. For HTTP-based
integration, start the server with \code{--transport http --port 8080}.

\subsection{Hardware requirements for data processing}
\label{app:hardware-requirements}

Beyond LLM memory, users must budget RAM for dataset operations.
\proglang{Rust} processes data in compact columnar format via
\pkg{polars}. The primary memory consumers are design matrices
($8nk$ bytes for $n$ observations and $k$ predictors), covariance
matrices ($8k^2$ bytes), bootstrap samples ($8kB$ bytes for $B$
replications), and panel transformations (approximately equal in size
to the original dataset).

\begin{table}[htbp]
\centering
\begin{threeparttable}
\caption{RAM requirements for data processing}
\label{tab:hardware-data}
\small
\begin{tabular}{lrrrl}
\toprule
Rows & File size & Base RAM & Matrix ops & Notes \\
\midrule
1,000 & $<$1 MB & 50 MB & +20 MB & Standard laptop \\
10,000 & 1--10 MB & 100 MB & +100 MB & Standard laptop \\
100,000 & 10--100 MB & 500 MB & +500 MB & Consumer hardware \\
1,000,000 & 100 MB--1 GB & 2 GB & +2 GB & Workstation \\
10,000,000 & 1--10 GB & 10 GB & +8 GB & High-memory system \\
\bottomrule
\end{tabular}
\begin{tablenotes}[flushleft]
\item \emph{Note:} Peak memory during analysis. Assumes typical
econometric datasets with 10--50 columns.
\end{tablenotes}
\end{threeparttable}
\end{table}

For combined LLM and data processing, sum the requirements from
Tables~\ref{tab:models} and \ref{tab:hardware-data}. A typical
configuration analyzing 100,000 observations with a 7B model requires
approximately 8--10~GB total RAM. Table~\ref{tab:hardware-summary}
provides system recommendations for common deployment scenarios.

\begin{table}[htbp]
\centering
\begin{threeparttable}
\caption{System configurations by deployment scenario}
\label{tab:hardware-summary}
\small
\begin{tabular}{llrrr}
\toprule
Scenario & Model & RAM & GPU VRAM & Storage \\
\midrule
Minimal (laptop) & Ministral 3B & 16 GB & 6 GB & 5 GB \\
Consumer (desktop) & Llama 4 Scout & 16 GB & 12 GB & 15 GB \\
Professional & Llama 3.3 70B & 64 GB & 48 GB & 50 GB \\
Server (large data) & Cloud API & 32 GB & --- & 2 GB \\
\bottomrule
\end{tabular}
\begin{tablenotes}[flushleft]
\item \emph{Note:} RAM includes both LLM and data processing headroom.
GPU requirements assume CUDA-capable NVIDIA cards; AMD and Intel GPUs
are supported via Ollama's ROCm and OneAPI backends.
\end{tablenotes}
\end{threeparttable}
\end{table}

\noindent For the simplest setup, the Docker configuration
(Appendix~\ref{app:deployment}) bundles all dependencies including
a pre-downloaded model and requires no local toolchain installation.

\section{Privacy and security} \label{app:privacy}
\setcounter{table}{0}
\setcounter{figure}{0}

This appendix documents data transmission behavior and security considerations.

\subsection{Data transmission summary}

Table~\ref{tab:data-transmission} lists operations that may transmit data
beyond the local machine when using cloud LLM providers.

\begin{table}[htbp]
\centering
\begin{threeparttable}
\caption{Data transmission by operation type}
\label{tab:data-transmission}
\small
\begin{tabular}{lp{6cm}l}
\toprule
Operation & Data Transmitted & Risk Level \\
\midrule
Dataset load & Column names, types, row count, null percentages & Low \\
Data preview (head/tail/sample) & Actual row values & High \\
Data profiling & Frequent values, min/max, examples & Medium \\
Statistical estimation & Coefficients, SEs, test statistics only & Low \\
Outlier detection & Specific observation values & Medium \\
Error messages & May include problematic values & Low \\
\bottomrule
\end{tabular}
\begin{tablenotes}[flushleft]
\item \emph{Note:} Applies only when using cloud LLM providers. Local
deployment via Ollama eliminates all external data transmission.
\end{tablenotes}
\end{threeparttable}
\end{table}

\subsection{Audit logging}

To enable audit logging of all MCP tool calls:

\begin{verbatim}
p2a-mcp --transport stdio --log-level debug --log-file audit.log
\end{verbatim}

\noindent The log records: timestamp, tool name, parameters (excluding bulk
data), and result summary. Review logs to verify no unexpected data transmission.

\subsection{Network verification}

To verify local-only operation with Ollama:

\begin{enumerate}
  \item Start Ollama: \code{ollama serve}
  \item Block outbound connections: \code{sudo ufw default deny outgoing}
  \item Run analysis---it should complete successfully
  \item Restore network: \code{sudo ufw default allow outgoing}
\end{enumerate}

\subsection{Security best practices}

\begin{itemize}
  \item Avoid loading untrusted data files (potential injection payloads)
  \item Review tool calls when processing sensitive data
  \item Use local LLM deployment for high-security environments
  \item Keep software updated; security patches are released as needed
  \item The MCP server has no filesystem access beyond loaded datasets
\end{itemize}

\noindent The primary mitigation for data privacy concerns is local
LLM deployment via Ollama, which eliminates all external data
transmission while maintaining full analytical capability.

\section{Reproducibility} \label{app:reproducibility}
\setcounter{table}{0}
\setcounter{figure}{0}

This appendix provides detailed instructions for reproducing the results in this paper.

\subsection{Version information}

Results in this paper were generated with the following software versions:
\begin{itemize}
  \item \proglang{Rust}: 1.92.0 (stable)
  \item \proglang{R}: 4.5.2 with OpenBLAS
  \item Key \proglang{R} packages: \pkg{plm} 2.6-4, \pkg{lfe} 3.1-0,
    \pkg{sandwich} 3.1-1, \pkg{fixest} 0.12.1, \pkg{forecast} 8.23.0
  \item Operating system: Ubuntu 24.04 LTS (Linux 6.8)
\end{itemize}

\subsection{Reproducing benchmarks}

Performance benchmarks can be reproduced using the following protocol:
\begin{enumerate}
  \item Install \proglang{R} 4.5+ with required packages: \pkg{plm}, \pkg{lfe},
    \pkg{sandwich}, \pkg{forecast}, \pkg{bench}
  \item Build \pkg{prompt2analytics} in release mode: \code{cargo build --release -p p2a-cli}
  \item Execute benchmarks: \code{make -C paper benchmark-full}
\end{enumerate}

\noindent Both languages use seed~$= 42$ for random number generation. The
benchmark protocol includes 10 warmup iterations before 100 measured iterations
for fast methods (50 for moderate, 20 for slow). Detailed
instructions: \code{paper/\allowbreak code/\allowbreak README\_REPRODUCIBILITY.md}.

\subsection{Containerized reproduction}

The repository's \code{docker-compose.yml} provides a containerized
deployment of the MCP server. For reproducing the paper's tables and
figures, the native toolchain described above is recommended; a
containerized paper build requires a multi-stage image with both
\proglang{R}~4.5.2 (plus packages listed above) and
\proglang{Rust}~1.92.0. A sample \code{Dockerfile} for this purpose
is provided in \code{paper/\allowbreak Dockerfile.reproduce}; after
building, run:

\begin{verbatim}
docker build -f paper/Dockerfile.reproduce -t p2a-paper .
docker run -v $(pwd)/paper:/output p2a-paper make all
\end{verbatim}

\noindent Results should match within the tolerances specified in
Appendix~\ref{app:benchmark-methodology}.

\section{Benchmark methodology} \label{app:benchmark-methodology}
\setcounter{table}{0}
\setcounter{figure}{0}

This appendix provides complete details on the benchmark methodology used in
Section~\ref{sec:deployment}.

\subsection{Data generating processes}

Benchmark datasets are generated synthetically using identical data-generating
processes in both languages with fixed random seeds (seed $= 42$) for exact
reproducibility.

For regression benchmarks, data follow
$y_i = \beta_0 + \beta_1 x_{1i} + \beta_2 x_{2i} + \varepsilon_i$ with
true parameters $\beta_0 = 1.0$, $\beta_1 = 2.0$, $\beta_2 = -1.5$, and
$\varepsilon_i \sim \mathcal{N}(0, \sigma^2)$ where $\sigma = 1.0$; covariates
are drawn as $x_{1i} \sim \mathcal{N}(0, 1)$ and $x_{2i} \sim \mathcal{N}(0, 1)$
with correlation $\rho = 0.3$. For panel data, entity effects
$\alpha_i \sim \mathcal{N}(0, 0.5^2)$ are added, with 100 entities observed
over 10 periods. For discrete choice, the latent index
$y^*_i = \beta' x_i + \varepsilon_i$ is transformed via the logit or probit
link. For time series, data follow ARIMA$(1,1,1)$ with $\phi = 0.7$,
$\theta = 0.3$, and innovation variance $\sigma^2_\eta = 1.0$.

Complete DGP specifications are provided in
\code{paper/code/README\_REPRODUCIBILITY\allowbreak.md}.

\subsection{Benchmark execution protocol}

Each benchmark configuration executes multiple iterations after a warmup phase:
\begin{itemize}
  \item 100 iterations for fast methods (OLS, robust SEs, PCA)
  \item 50 iterations for moderate methods (panel FE, clustering)
  \item 20 iterations for slow methods (ARIMA, MSTL (multiple seasonal-trend
    decomposition using LOESS))
\end{itemize}

\proglang{R} benchmarks use the \pkg{bench} package, which accounts for garbage
collection timing. \proglang{Rust} benchmarks use a custom harness
(\code{harness = false} in \code{Cargo.toml}) with automatic warmup
and outlier detection. We report full timing distributions (minimum,
quartiles, maximum, mean, standard deviation) rather than point estimates alone.

\subsection{Hardware and software environment}

All benchmarks run on identical hardware: an Intel Core i7-1260P processor
with 64~GB RAM running Linux, with \proglang{R} 4.5.2 linked against
OpenBLAS and \proglang{Rust} 1.92.0 compiled in release mode with
link-time optimization. The \pkg{faer} library (version~0.22) used by
\pkg{p2a-core} does \emph{not} call an external BLAS; it implements its
own GEMM microkernels in pure \proglang{Rust} with explicit SIMD via
the \pkg{pulp} crate, dispatching at runtime to AVX2+FMA on x86-64
(or NEON on ARM). Both \proglang{R} and \proglang{Rust} therefore
use optimized linear algebra routines, but from different
implementations (OpenBLAS vs.\ \pkg{faer}). This makes the
comparison a realistic measure of what each ecosystem delivers out of
the box, not a controlled BLAS-versus-BLAS experiment.

\subsection{Validation tolerances} \label{app:validation-tolerances}

Numerical validation against \proglang{R} reference implementations uses
sample-size-dependent tolerances. We use absolute tolerances for quantities
with known scale (e.g., test statistics) and relative tolerances for
unbounded quantities (e.g., coefficient estimates).

\begin{table}[htbp]
\centering
\begin{threeparttable}
\caption{Validation tolerances by method and sample size}
\label{tab:validation-tolerances}
\small
\begin{tabular}{llccc}
\toprule
Category & Quantity & Small ($n < 100$) & Medium & Large ($n > 10{,}000$) \\
\midrule
\multicolumn{5}{l}{\textit{Coefficients and standard errors}} \\
OLS coefficients & Relative & $10^{-6}$ & $10^{-8}$ & $10^{-10}$ \\
Robust SEs (HC0/HC1)$^a$ & Relative & $10^{-6}$ & $10^{-8}$ & $10^{-8}$ \\
Robust SEs (HC2/HC3)$^a$ & Relative & $10^{-5}$ & $10^{-7}$ & $10^{-7}$ \\
Panel FE/RE & Relative & $10^{-5}$ & $10^{-6}$ & $10^{-6}$ \\
MLE$^b$ (logit/probit) & Relative & $10^{-4}$ & $10^{-6}$ & $10^{-6}$ \\
\addlinespace
\multicolumn{5}{l}{\textit{Test statistics and fit measures}} \\
$t$-statistics, $F$-statistics & Absolute & $10^{-4}$ & $10^{-6}$ & $10^{-6}$ \\
$R^2$, within $R^2$ & Absolute & $10^{-4}$ & $10^{-6}$ & $10^{-6}$ \\
$p$-values & Absolute & $10^{-4}$ & $10^{-6}$ & $10^{-6}$ \\
Log-likelihood & Absolute & $10^{-2}$ & $10^{-4}$ & $10^{-4}$ \\
\addlinespace
\multicolumn{5}{l}{\textit{Hypothesis tests}} \\
Test statistics & Absolute & $10^{-4}$ & $10^{-6}$ & $10^{-6}$ \\
Critical values & Absolute & $10^{-4}$ & $10^{-4}$ & $10^{-4}$ \\
\bottomrule
\end{tabular}
\begin{tablenotes}[flushleft]
\item \emph{Note:} Absolute tolerances ($\varepsilon_{\mathrm{abs}}$) apply to
test statistics, $p$-values, and bounded quantities. Relative tolerances
($\varepsilon_{\mathrm{rel}}$) apply to coefficient estimates and standard errors.
$^a$~HC0--HC3: heteroskedasticity-consistent standard error estimators
(types 0--3). $^b$~MLE: maximum likelihood estimation.
\end{tablenotes}
\end{threeparttable}
\end{table}

Looser tolerances for small samples reflect
increased sensitivity to algorithmic differences when matrices are
ill-conditioned. HC2/HC3 standard errors use looser tolerances because
leverage calculation involves $h_{ii} = x_i'(X'X)^{-1}x_i$, which accumulates
rounding errors differently depending on whether leverage is extracted from
QR decomposition (as in \pkg{sandwich}) or computed via direct matrix
multiplication (as in \pkg{prompt2analytics}). Both approaches are numerically
valid; the tolerance reflects implementation variation rather than error.

MLE methods use looser small-sample tolerances because convergence criteria
differ between implementations: \proglang{R}'s \fct{glm} uses relative
change in deviance, while \pkg{prompt2analytics} uses absolute change in
log-likelihood. Final parameter estimates can differ at the tolerance boundary
while both satisfying their respective convergence criteria.

Tolerances were determined empirically by
running validation tests across sample sizes from $n = 10$ to $n = 100{,}000$
and identifying the tightest tolerance that passes consistently. Edge cases
(singular matrices, near-separation in logit/probit, small cluster counts)
are tested separately with method-specific tolerances documented in the
test code. All 416 validation tests pass at the specified tolerances.


%% Appendix I (Related work) folded into Section 2.3

\section[Numerical validation against R]{Numerical validation against \proglang{R}} \label{app:numerical-validation}
\setcounter{table}{0}
\setcounter{figure}{0}

We validate all methods against reference \proglang{R} implementations using
sample-size-dependent tolerances (looser for small samples where
ill-conditioning amplifies differences, tighter for large samples).
Reference packages include \fct{lm}/\pkg{sandwich} for OLS,
\pkg{plm}/\pkg{lfe} for panel data, \fct{glm} for discrete choice,
and \pkg{forecast} for time series. All validation tests use fixed
random seeds (seed $= 42$) for reproducibility; reference \proglang{R}
scripts are provided in the supplementary materials.

The Longley dataset~\citep{Longley:1967}, with its extreme
multicollinearity (condition number $> 10^9$), provides a stringent
test. Table~\ref{tab:validation-ols} shows coefficients match
\proglang{R}'s \fct{lm} to at least 8 significant digits.

\begin{table}[htbp]
\centering
\begin{threeparttable}
\caption{OLS validation: Longley dataset}
\label{tab:validation-ols}
\renewcommand{\arraystretch}{1.15}
\begin{tabular}{lrrr}
\toprule
Coefficient & \proglang{R} \fct{lm} & \pkg{prompt2analytics} & Difference \\
\midrule
(Intercept) & $-3482.259$ & $-3482.259$ & $< 10^{-8}$ \\
GNP.deflator & $0.01506$ & $0.01506$ & $< 10^{-8}$ \\
GNP & $-0.03582$ & $-0.03582$ & $< 10^{-8}$ \\
Unemployed & $-0.02020$ & $-0.02020$ & $< 10^{-8}$ \\
Armed.Forces & $-0.01033$ & $-0.01033$ & $< 10^{-8}$ \\
Population & $-0.05110$ & $-0.05110$ & $< 10^{-8}$ \\
Year & $1.82915$ & $1.82915$ & $< 10^{-8}$ \\
\midrule
$R^2$ & $0.9955$ & $0.9955$ & $< 10^{-4}$ \\
\bottomrule
\end{tabular}
\begin{tablenotes}[flushleft]
\item \emph{Note:} $n = 16$, $k = 6$. The Longley dataset's extreme
multicollinearity (condition number $> 10^9$) provides a stringent
near-singular test case: the design matrix is close to rank-deficient,
exercising the Cholesky-based solver's tolerance framework. All
coefficients match to at least 8 significant digits, confirming that
the implementation handles ill-conditioned problems correctly.
\end{tablenotes}
\end{threeparttable}
\end{table}

Robust standard errors (HC0--HC3) match
\pkg{sandwich}~\citep{Zeileis:2004} to $< 10^{-7}$; panel fixed
effects match \pkg{plm}~\citep{Croissant+Millo:2008} to $< 10^{-6}$
on the Grunfeld dataset; two-way HDFE via the MAP
algorithm~\citep{Gaure:2013} matches \pkg{lfe}'s \fct{felm} to
$< 10^{-5}$.

Beyond the methods shown above, the validation suite covers over 200
methods with more than 1,800 individual test cases (437 specifically
targeting validation against \proglang{R} reference implementations or
known data-generating processes).
Statistical tests ($t$-tests, ANOVA, chi-squared, Fisher's exact,
Wilcoxon, Shapiro-Wilk, Kolmogorov-Smirnov) are validated against
\proglang{R}'s \pkg{stats} package. Discrete choice models (logit,
probit) match \fct{glm} coefficients and standard errors. Time series
methods (VAR, ARIMA) match \pkg{vars} and \pkg{forecast}. The \pkg{p2a-core} test suite comprises 2{,}168 test functions,
of which 442 are validation tests that compare output against
\proglang{R} reference implementations or known data-generating
processes (Table~\ref{tab:test-distribution}).  Every module with
user-facing methods has at least 10 tests, and the most complex
modules (panel, matching, TMLE, HDFE) have 17--26 each.
Continuous integration (GitHub Actions) runs the full test suite,
\code{clippy} linting, and format checks on every push to
\code{main} and on all pull requests. Line-level code coverage is not
currently tracked as a CI gate; running \code{cargo-tarpaulin} on the
full suite yields approximately 62\% line coverage for \pkg{p2a-core},
reflecting the large number of edge-case code paths (error handling,
fallback solvers, GPU dispatch stubs) that are exercised only under
specific hardware or data conditions.  We consider validation breadth
(442 cross-validated tests spanning all method categories) a more
informative quality signal than aggregate line coverage for a
numerical library.  The CI pipeline does \emph{not} execute the
\proglang{R} reference scripts; cross-language validation is run
manually before releases and the reference outputs are checked into
the repository.  All Rust-side validation tests can be run via
\code{cargo test -p p2a-core}.

\paragraph{Monte Carlo validation of statistical properties.}
Point-estimate matching against \proglang{R} verifies that the same
numbers are produced, but does not test whether confidence intervals,
standard errors, and $p$-values have correct statistical properties.
A separate Monte Carlo validation framework
(\code{paper/code/mc\_validation/}) addresses this by running repeated
simulations under known data-generating processes and checking five
properties: (i)~95\% confidence interval coverage rates, (ii)~standard
error calibration (ratio of mean reported SE to empirical SD of
estimates), (iii)~Type~I error rates at $\alpha = 0.05$, (iv)~negative
controls confirming that standard SEs \emph{fail} under
heteroskedasticity, and (v)~statistical power under the alternative.
Tolerances are derived from binomial confidence intervals around the
nominal rates, accounting for finite-simulation variability.

Across 29 methods (OLS with HC0--HC3, panel FE/RE, logit, probit,
IV/2SLS, DiD, IPW, TMLE, and nine hypothesis tests) and 94 test
scenarios at sample sizes $n \in \{200, 1{,}000\}$ with 2{,}000
Monte Carlo replications each, 92 pass (97.9\%).  Two additional
negative-control scenarios (standard OLS SEs applied to
heteroskedastic data) correctly detect under-coverage, confirming
that robust SEs are needed.  Standard error calibration and Type~I
error control are exact: all 34 SE accuracy tests pass (mean ratio
1.001) and all 16 Type~I error tests pass.  Mean confidence interval
coverage is 0.950 (nominal 0.950).  The two substantive failures are
marginal over-coverage at $n = 1{,}000$ (OLS 96.0\%, logit 96.3\%),
both within one percentage point of the nominal rate.  IPW uses the
efficient influence function variance estimator of
\citet{Lunceford2004}, which accounts for propensity score estimation
uncertainty, achieving 94.3--94.9\% coverage (SE ratio $\approx
1.0$).  Full results and the simulation code are included in the
supplementary materials.

\begin{table}[htbp]
\centering
\begin{threeparttable}
\caption{Test distribution by module category}
\label{tab:test-distribution}
\small
\begin{tabular}{lrrr}
\toprule
Module & Tests & Validation$^a$ & Methods \\
\midrule
Regression (OLS, diagnostics, NLS, LOESS, GLS) & 120 & 55 & 25 \\
Panel data (FE, RE, HDFE, GMM) & 95 & 40 & 12 \\
Causal inference (DiD, IV, RD, synth) & 130 & 52 & 20 \\
Treatment effects (IPW, TMLE, DML, matching) & 145 & 60 & 18 \\
Discrete choice (logit, probit, count) & 55 & 22 & 9 \\
Time series (VAR, ARIMA, GARCH, Kalman) & 110 & 45 & 20 \\
Statistical tests & 160 & 65 & 50+ \\
ML (clustering, PCA, t-SNE, SVM) & 120 & 35 & 10 \\
Spatial econometrics & 50 & 18 & 8 \\
Survival analysis & 35 & 12 & 5 \\
Data, simulation, export, visualization & 148 & 38 & --- \\
\midrule
\textbf{Total} & \textbf{2{,}168} & \textbf{442} & \textbf{250+} \\
\bottomrule
\end{tabular}
\begin{tablenotes}[flushleft]
\item $^a$~Tests that compare against \proglang{R} reference
implementations or verify recovery of known parameters from simulated
data-generating processes.
\end{tablenotes}
\end{threeparttable}
\end{table}

For complex causal inference and treatment effect methods, validation
depth varies. Canonical 2$\times$2 difference-in-differences, IV/2SLS,
IPW, doubly robust estimation (AIPW), and causal mediation analysis
are cross-validated against \proglang{R} packages---\fct{lm} for DiD,
\pkg{AER} for IV, manual Horvitz--Thompson/H\'{a}jek estimators for
IPW, and the \pkg{mediation} package---on
synthetic datasets with known treatment effects (e.g., true ATT~$=5$
for DiD, true $\beta=1$ for IV, true ATE~$=0.75$ for IPW/AIPW).
Propensity score matching is validated against
\pkg{MatchIt}~\citep{Ho+etal:2011} for balance diagnostics
(standardized mean differences, variance ratios).  Regression
discontinuity is validated against the \pkg{rdrobust}
framework~\citep{Calonico+etal:2014}, checking that local
linear estimates recover the true discontinuity within tolerance.
FEGLM is validated against \pkg{fixest}'s
\fct{feglm}~\citep{Berge:2018} for coefficient sign, pseudo~$R^2$,
and the logit/probit scaling ratio~($\approx 1.6$).  TMLE has eight
validation tests that verify the targeting step, clever covariate
formula, influence curve properties, propensity score truncation, and
counterfactual predictions on simulated data; these check mathematical
identities and DGP-implied ranges rather than exact \proglang{R}
output values. In total, 41 \proglang{R} validation scripts and
73 expected-value CSVs are provided in the supplementary materials.

Several complex methods---staggered
DiD~\citep{Callaway+SantAnna:2021}, extended TWFE, Goodman-Bacon
decomposition, Double/Debiased ML, synthetic control, SCPI, CTMLE,
LTMLE, CBPS, and entropy balancing---are validated using
DGP-based tests with known true parameters rather than exact numerical
cross-validation against \proglang{R} packages.  These tests generate
data from a known process (e.g., true ATT~$=5$, true coefficient~$=1$)
and verify that the estimator recovers the parameter within tolerance,
that pre-trends are null when they should be, that weights sum correctly,
and that confidence intervals have appropriate coverage.
The corresponding \proglang{R} implementations (e.g.,
\pkg{did}, \pkg{DoubleML}, \pkg{Synth}) use different default tuning
parameters, cross-validation splits, and optimization algorithms that
complicate exact numerical comparison.  Extending
exact cross-validation coverage to these methods is a priority for future
releases.

Appendix~\ref{app:benchmark-methodology} documents tolerance
thresholds by method category and sample size.
Appendix~\ref{app:implementation} lists known differences from
reference implementations (default robust SE type, convergence
tolerances, $R^2$ variant selection)---none affecting the validity
of statistical inference.

For core panel methods, \pkg{prompt2analytics} provides comparable
coverage to specialized \proglang{R} packages; the primary remaining
gap is Conley spatial standard errors. Users requiring maximum
performance for high-dimensional fixed effects should note that
\pkg{fixest}~\citep{Berge:2018} likely outperforms
\pkg{prompt2analytics} due to its highly optimized C++
implementation. The distinguishing features are the natural language
interface and fully local execution without \proglang{R} dependencies;
the package complements rather than replaces specialized tools.

\section{Implementation notes} \label{app:implementation}
\setcounter{table}{0}
\setcounter{figure}{0}

This appendix documents implementation caveats and limitations for specific
methods.

\subsection{Instrumental variables}

The two-stage least squares (2SLS) implementation reports first-stage
$F$-statistics for instrument relevance. The traditional rule of thumb ($F > 10$) from \citet{Staiger+Stock:1997}
has been superseded by more nuanced guidance. \citet{Stock+Yogo:2005} provide
critical values depending on the number of instruments and acceptable bias
levels; \citet{Lee+etal:2022} offer further refinements for inference under
weak identification. The current implementation does not provide Stock-Yogo
critical values or weak-identification-robust inference methods (Anderson-Rubin,
conditional likelihood ratio). Users with potentially weak instruments should
verify results using \proglang{R}'s \pkg{ivmodel} or \pkg{AER} packages.

\subsection{Difference-in-differences}

The difference-in-differences (DiD) implementation assumes the canonical
two-group, two-period design.
Recent econometric literature has documented severe biases when this estimator
is applied to staggered treatment timing. \citet{Goodman-Bacon:2021} demonstrates
that two-way fixed effects DiD produces a weighted average of all possible $2 \times 2$
comparisons, including problematic ``forbidden'' comparisons using already-treated
units as controls. \citet{deChaisemartin+DHaultfoeuille:2020} and
\citet{Sun+Abraham:2021} propose alternative estimators robust to treatment
effect heterogeneity.

\pkg{prompt2analytics} implements staggered DiD via
Callaway-Sant'Anna~\citep{Callaway+SantAnna:2021}, Goodman-Bacon decomposition, and extended
two-way fixed effects (TWFE)~\citep{Wooldridge:2025}. Users with complex staggered designs may also
consult \pkg{fixest}'s \fct{sunab} function or the \pkg{did} package in
\proglang{R} for additional diagnostics.

\subsection{Random effects}

The random effects estimator uses the Swamy-Arora variance component estimator,
which can produce negative variance estimates for $\sigma^2_\alpha$ (the
between-entity variance) in certain data configurations. When this occurs,
the implementation truncates the estimate to zero and proceeds with pooled
OLS, issuing a warning.

\subsection{Discrete choice models}

Logit and probit models report coefficients on the latent index scale. For
interpretation, marginal effects are computed as \emph{Average Marginal Effects}
(AME): the partial derivative $\partial \Pr(y=1)/\partial x_j$ is evaluated
at each observation's covariate values, then averaged across the sample.
This differs from \emph{Marginal Effects at Means} (MEM), which evaluates
derivatives at the sample mean of all covariates. AME is generally preferred
because it does not depend on potentially unrepresentative ``average''
covariate values~\citep{Cameron+Trivedi:2005}.

\subsection{Time series models}

Autoregressive integrated moving average (ARIMA) estimation supports
non-seasonal models: ARIMA$(p,d,q)$ with automatic or user-specified orders. Seasonal ARIMA (SARIMA) is partially implemented:
seasonal differencing and seasonal AR/MA terms can be specified, but automatic
seasonal order selection is not available.

Vector autoregression (VAR) and vector error correction model (VECM)
implementations support lag selection via AIC/BIC and Johansen
cointegration testing, but do not implement structural VAR identification.

\subsection{Regression diagnostics validation}

Diagnostic tests are validated against \pkg{lmtest}~\citep{lmtest}:
Jarque-Bera normality test statistics match to $< 10^{-4}$,
Breusch-Pagan heteroskedasticity tests to $< 10^{-4}$,
Durbin-Watson autocorrelation statistics to $< 10^{-6}$, and
variance inflation factors (VIF) to $< 10^{-6}$.

\subsection{Known differences from reference implementations}
\label{app:known-differences}

A small number of intentional differences exist between \pkg{prompt2analytics}
and reference implementations. First, \pkg{prompt2analytics} defaults to
HC1 robust standard errors (consistent with \proglang{Stata}'s \code{robust} option),
while \proglang{R}'s \pkg{sandwich} defaults to HC3. Second, \proglang{R}
uses variable names in output; \pkg{prompt2analytics} uses numeric indices
internally but displays names in formatted output. Third, for panel models
we report ``within $R^2$'' by default, matching \pkg{plm}'s convention for
fixed effects models. Fourth, iterative methods (HDFE, maximum likelihood estimation) use a default
convergence tolerance of $10^{-8}$, which may differ slightly from
\proglang{R} package defaults. These differences are documented in the
method documentation and do not affect the validity of statistical inference.

\noindent Coverage gaps and unimplemented methods are documented in
Appendix~\ref{app:best-practices} (Table~\ref{tab:method-coverage}).

\section{Extended examples} \label{app:extended-examples}
\setcounter{table}{0}
\setcounter{figure}{0}

This appendix collects additional examples that complement the core
demonstrations in Section~\ref{sec:examples}.

\subsection{Command-line interface} \label{app:cli-examples}

The command-line interface provides programmatic access for scripted
workflows.\footnote{The output shown here was captured from
actual CLI invocations; the full transcript is in
\code{paper/transcripts/section4\_2\_cli\_output.txt}.}
Commands follow a hierarchical structure with sessions maintaining state
across invocations:

\begin{Sinput}
$ p2a --session analysis.json data load grunfeld.csv --name grunfeld
\end{Sinput}

\begin{Soutput}
Dataset: grunfeld
  Rows: 200
  Columns: 5
  Column names: firm, year, inv, value, capital
\end{Soutput}

\begin{Sinput}
$ p2a --session analysis.json reg ols grunfeld \
    -y inv -x value capital --robust hc1
\end{Sinput}

\begin{Soutput}
OLS Regression Results
============================================================
Variable         Coefficient   Std. Error    t-value    p-value
------------------------------------------------------------
(Intercept)       -42.714369    11.574701     -3.690     0.0003***
value               0.115562     0.006811     16.967     0.0000***
capital             0.230678     0.048866      4.721     0.0000***
------------------------------------------------------------
R-squared: 0.812408
Adj. R-squared: 0.809537
Observations: 200
\end{Soutput}

\begin{Sinput}
$ p2a --session analysis.json panel fe grunfeld \
    -y inv -x value capital --entity firm
\end{Sinput}

\begin{Soutput}
Fixed Effects (Within) Estimation Results
============================================================
Variable         Coefficient   Std. Error    t-value    p-value
------------------------------------------------------------
value               0.110124     0.011857      9.288     0.0000***
capital             0.310065     0.017355     17.867     0.0000***
------------------------------------------------------------
R-squared: 0.766758,  Observations: 200,  Groups: 10
F-statistic: 309.0142 (p-value: 0.0000)
\end{Soutput}

\noindent The session mechanism records all commands and can export them
as an executable shell script for exact reproduction.

\subsection[Comparison with R]{Comparison with \proglang{R}}

For reference, the equivalent \proglang{R} code for the fixed effects
estimation shown in Section~\ref{sec:examples-chat} is:

\begin{Sinput}
library(plm)
data("Grunfeld", package = "plm")
fe_model <- plm(inv ~ value + capital, data = Grunfeld,
                index = c("firm", "year"), model = "within")
summary(fe_model)
\end{Sinput}

The \pkg{prompt2analytics} implementation produces numerically equivalent
results (Section~\ref{sec:evaluation}), while offering a conversational
interface for users unfamiliar with \proglang{R} syntax.

\subsection{Time series analysis}

The basic time series workflow (stationarity testing and ARIMA fitting)
is demonstrated in Section~\ref{sec:examples-ts}. Here we note a
limitation: order selection is delegated to the LLM rather than
algorithmic search, so it is not equivalent to \proglang{R}'s
\code{auto.arima()}~\citep{Hyndman+Khandakar:2008}.

\subsection{Handling messy data}

Real-world datasets often require cleaning before analysis. The
system detects non-standard missing indicators and guides type
conversion:

\begin{Sinput}
User: Load survey_responses.csv and describe the data.
\end{Sinput}

\begin{Soutput}
Dataset: survey_responses (2,847 rows x 6 columns)
Note: 'income' and 'satisfaction' are typed as String, likely due to
non-standard missing value indicators (".", "N/A").
\end{Soutput}

\begin{Sinput}
User: Replace "." and "N/A" with proper missing values, then convert
      income and satisfaction to numeric.
\end{Sinput}

\begin{Soutput}
Cleaned: 89 values replaced across 2 columns.
  income:        now Float64, 52 missing (1.8%)
  satisfaction:  now Float64, 37 missing (1.3%)
\end{Soutput}

\noindent The next example demonstrates a common challenge: ambiguous
column names that do not match natural language descriptions. The
dataset uses coded variable names rather than descriptive labels.

\begin{Sinput}
User: Load hr_extract_q3.csv and describe the data.
\end{Sinput}

\begin{Soutput}
Dataset: hr_extract_q3 (1,204 rows x 8 columns)
  emp_id:    Int64 (1204 unique)
  dept_cd:   String (12 unique: "ENG", "MKT", "FIN", ...)
  var_03:    Float64 (range: 28000-185000)
  var_04:    Float64 (range: 0.5-35.2)
  var_05:    Int64 (0 or 1)
  q3_rating: Float64 (range: 1.0-5.0, 23 missing)
  loc_cd:    String (8 unique: "NYC", "SFO", "CHI", ...)
  dte_hire:  String (formats vary: "2015-03-12", "03/12/2015")
\end{Soutput}

\begin{Sinput}
User: var_03 is salary and var_04 is years of experience.
      Regress salary on experience and department with robust SEs.
\end{Sinput}

\begin{Soutput}
Renamed: var_03 -> salary, var_04 -> experience
Note: dept_cd is categorical (12 levels); using treatment coding
with reference category "ENG".

OLS Regression (HC1 robust standard errors)
  Dependent variable: salary
  Observations: 1,204
  ...
\end{Soutput}

\noindent This example illustrates both a strength and a limitation.
The system correctly maps user-provided variable descriptions to coded
column names and handles the categorical department variable
automatically. However, the user must know what the coded variables
represent---the system cannot infer that \code{var\_03} is salary
from its range alone.  Inconsistent date formats in \code{dte\_hire}
(ISO~8601 and US~format in the same column) would require explicit
cleaning before use as a time variable.
